\uuid{aYol}
\titre{Primitive}
\auteur{bodin ; croizat ; sacre}
\organisation{exo7}
\licence{CC BY-NC-SA 4.0}
\datecreate{2022-03-08}
\niveau{L1}
\module{Analyse}
\chapitre{Calcul d'intégrales}
\sousChapitre{Primitives diverses}
\difficulte{3}
\theme{Primitive}
\exo7qcmid{hilisit.equations-differentielles.011}

\contenu{
\begin{question}
\qclassification{127.11}
Quelles sont les affirmations vraies ?
\begin{answers}
  \good{Une primitive de $\sin(x)e^{\cos(x)}$ est $-e^{\cos(x)}$.}
  \bad{Une primitive de $\cos(x^3+x)$ est $\sin(x^3+x)$.}
  \good{Une primitive de $\ln(x)$ est $x\ln(x)-x$ (sur $]0,+\infty[$).}
  \good{Une primitive de $4x^3+4x$ est $(x^2+1)^2$.}
\end{answers}
\begin{explanations}
Une primitive de $\sin(x)e^{\cos(x)}$ est bien $-e^{\cos(x)}$ puisque $(-e^{\cos(x)})' = -(-\sin(x))e^{\cos(x)} = \sin(x) e^{\cos(x)}$.
La dérivée de $\sin(x^3+x)$ est $(3x^2+1)\cos(x^3+x)$, donc $\sin(x^3+x)$ n'est pas une primitive de $\cos(x^3+x)$. Oui une primitive de $\ln(x)$ est $x\ln(x)-x$ puisque la dérivée de cette-dernière donne bien $\ln(x) + x \cdot \frac 1x - 1 = \ln(x)$.
Enfin, la dérivée de $(x^2+1)^2$ est $2 \times 2x \times (x^2+1) = 4x^3+4x$, donc  $(x^2+1)^2$ est bien une primitive de $4x^3+4x$.
\end{explanations}
\end{question}
}
